% ============================================================
%  Chapter: State of the Art
%  Multi-Party Speech-Based Conversational AI
% ============================================================

\section{State of the Art}
\label{sec:sota}

This chapter surveys the research landscape relevant to the design of
speech-based conversational agents that participate in multi-party group
interactions.  The discussion is organised along four thematic axes.
Section~\ref{subsec:speech-architectures} reviews the architectural
families for building speech-based conversational systems---cascaded,
half-cascade, and end-to-end---and compares their trade-offs.
Section~\ref{sec:turntaking} examines turn-taking mechanisms, from
foundational linguistic models to recent machine-learning approaches,
with particular attention to the challenges introduced by multi-party
settings.  Section~\ref{sec:vad-vap} covers Voice Activity Detection and
Voice Activity Projection, the low-level perceptual building blocks that
underpin turn management.  Section~\ref{subsec:multiparty-ca} presents
multi-party conversational AI as an integrating theme, reviewing
taxonomies, existing systems, and technical challenges.  Finally,
Section~\ref{subsec:gaps} identifies eight research gaps at the
intersection of these threads, and Section~\ref{subsec:positioning}
discusses the positioning of AIutami within the resulting landscape.

% --- Speech-Based Conversational Architectures ---
% ============================================================
%  \subsection: Speech-Based Conversational Architectures
%  Parent: \section{State of the Art}
% ============================================================

\subsection{Speech-Based Conversational Architectures}
\label{subsec:speech-architectures}

The design of speech-based conversational agents revolves around a fundamental
question: how should a system map an incoming acoustic signal to a meaningful
spoken response?  Ji et al.\ provide a comprehensive taxonomy in the WavChat
survey~\cite{ji2024wavechat}, identifying three broad families of
architectures---\emph{cascaded}, \emph{half-cascade} (or semi-cascaded), and
\emph{end-to-end}---that differ in the degree to which speech understanding,
language reasoning, and speech synthesis are fused into a single model.  This
subsection reviews each family and compares their trade-offs.

% ----------------------------------------------------------
\subsubsection{Cascaded Pipeline (STT--LLM--TTS)}
\label{sssec:cascaded}

The cascaded approach decomposes the problem into three independently developed
stages: an Automatic Speech Recognition (ASR) module transcribes the user's
utterance into text, a Large Language Model (LLM) produces a textual response,
and a Text-to-Speech (TTS) engine renders that response as audio.  Because each
component communicates through a well-defined textual interface, the pipeline is
highly modular: individual stages can be upgraded, replaced, or scaled without
retraining the others.

This architecture underpins the majority of commercially deployed voice
assistants, including the speech pipelines offered by Azure, Google Cloud, and
Amazon Web Services.  In the research literature, X-Talk~\cite{lin2025xtalk}
proposes a cascaded dialogue framework targeting low-latency interaction, while
ESPnet-SDS~\cite{li2025espnetsds} provides an open-source toolkit that
implements cascaded spoken-dialogue systems with pluggable ASR, LLM, and TTS
back-ends.

The principal advantages of the cascaded pipeline are its maturity and
interpretability.  The intermediate text representation enables straightforward
logging, debugging, and content moderation---capabilities that are essential in
enterprise and safety-critical settings.  Training costs are low because each
stage can leverage pre-trained, off-the-shelf models.

The main drawback is latency.  The sequential execution of three distinct
inference steps introduces cumulative delays in the range of two to five
seconds, far exceeding the approximately 200\,ms inter-turn gap that
characterises natural human conversation~\cite{heldner2010pauses}.  Moreover,
errors compound across stages: an ASR misrecognition propagates into the LLM
prompt and may trigger an irrelevant response.  Finally, the text bottleneck
discards paralinguistic information---prosody, emotion, hesitation
markers---that is often crucial for pragmatic interpretation, and prevents
native support for full-duplex interaction.

% ----------------------------------------------------------
\subsubsection{Half-Cascade (Semi-Cascaded) Architectures}
\label{sssec:half-cascade}

Half-cascade designs seek a middle ground by fusing two of the three stages
while retaining the third as a separate component.  The WavChat
taxonomy~\cite{ji2024wavechat} distinguishes three variants, each eliminating
a different inter-stage boundary.

\emph{Speech-In fusion} (Variant~A) removes the ASR stage: the LLM ingests
audio features directly through a learnable encoder and projection layer,
preserving paralinguistic cues that text transcription would discard (e.g.,
SALMONN~\cite{tang2024salmonn}, Freeze-Omni~\cite{wang2024freezeomni}).
\emph{Speech-Out fusion} (Variant~B) absorbs the TTS stage into the LLM,
which generates discrete speech tokens alongside text, reducing output latency
(e.g., LLaMA-Omni~\cite{fang2024llamaomni}).  \emph{Inner Monologue}
(Variant~C) retains text as an internal reasoning scaffold while producing
speech output directly, attempting to combine coherent reasoning with
expressive synthesis (e.g., SpeechGPT~\cite{zhang2023speechgpt}).

These architectures reduce latency and exploit richer multimodal
representations compared to cascaded pipelines.  However, they introduce
significant trade-offs.  Aligning audio and text representation spaces
requires carefully staged training curricula, and there is a documented risk
of \emph{reasoning degradation}---the model's language abilities may decline
after fine-tuning on speech data~\cite{chen2025urobench}.  Speech quality
also tends to lag behind dedicated TTS systems, and none of these models has
been evaluated beyond dyadic interaction.

% ----------------------------------------------------------
\subsubsection{End-to-End Speech-to-Speech Models}
\label{sssec:e2e}

End-to-end architectures collapse the entire spoken-dialogue pipeline into a
single neural model that maps audio input directly to audio output.  A key
enabler of this paradigm is the family of \emph{neural audio codecs}---models
such as SoundStream, EnCodec, and Mimi---that convert continuous waveforms into
compact sequences of discrete tokens amenable to autoregressive modelling.

Moshi~\cite{defossez2024moshi} is widely regarded as the first publicly
demonstrated full-duplex speech-to-speech model, achieving a mouth-to-ear
latency of approximately 200\,ms.  Commercial systems such as GPT-4o
(OpenAI), Gemini~2.5 (Google), Qwen3-Omni (Alibaba), and Kimi Audio
(Moonshot AI) have since adopted end-to-end or near-end-to-end designs, though
their architectural details remain largely proprietary.  In the academic domain,
dGSLM extends the Generative Spoken Language Model to dialogue by modelling two
audio channels jointly.

The appeal of end-to-end models is clear: by eliminating intermediate
representations, they minimise latency and preserve paralinguistic cues---pitch
contour, speaking rate, emotional colouring---that cascaded systems discard.
They also support full-duplex interaction natively, as the model can listen and
speak simultaneously.

Nevertheless, current end-to-end systems exhibit significant limitations.
URO-Bench~\cite{chen2025urobench} demonstrates that, when evaluated on
reasoning-intensive tasks, end-to-end speech models substantially underperform
their text-only LLM counterparts, suggesting that the speech modality introduces
a reasoning bottleneck.  Training costs are high, requiring large volumes of
paired speech--speech data.  Reliability remains a concern: Lin et
al.~\cite{lin2025xtalk} observe that end-to-end speech models are ``often
unsuitable for deployment'' in production settings due to hallucination,
inconsistent output quality, and limited controllability.

% ----------------------------------------------------------
\subsubsection{Comparative Summary}
\label{sssec:arch-comparison}

Table~\ref{tab:arch-comparison} synthesises the trade-offs discussed above
across nine evaluation dimensions.

\begin{table}[ht]
  \centering
  \caption{Comparison of speech-based conversational architectures.}
  \label{tab:arch-comparison}
  \small
  \begin{tabularx}{\textwidth}{@{} l >{\centering\arraybackslash}X >{\centering\arraybackslash}X >{\centering\arraybackslash}X @{}}
    \toprule
    \textbf{Aspect} & \textbf{Cascaded} & \textbf{Half-Cascade} & \textbf{End-to-End} \\
    \midrule
    Latency
      & High (2--5\,s)
      & Medium
      & Low (${\sim}$200\,ms) \\
    Reasoning quality
      & High (text LLM)
      & Medium (risk of degradation)
      & Low~\cite{chen2025urobench} \\
    Paralinguistic info
      & Lost (text bottleneck)
      & Partially preserved
      & Fully preserved \\
    Audio quality
      & High (dedicated TTS)
      & Medium
      & Variable \\
    Full-duplex
      & Not native
      & Partial
      & Native \\
    Modularity
      & High
      & Medium
      & Low \\
    Training cost
      & Low
      & Medium
      & High \\
    Production readiness
      & High
      & Medium
      & Low~\cite{lin2025xtalk} \\
    Multi-party support
      & Demonstrated
      & Unexplored
      & Unexplored \\
    \bottomrule
  \end{tabularx}
\end{table}

% ----------------------------------------------------------
\subsubsection{The Multi-Party Gap}
\label{sssec:multiparty-gap}

A striking observation emerges from the literature: virtually all end-to-end and
half-cascade models are designed and evaluated exclusively in \emph{dyadic}
settings, i.e., one human conversing with one agent.  Multi-party spoken
interaction---where an AI agent participates in a group conversation with two or
more human speakers---remains largely unexplored in the context of neural
speech-to-speech models.  Prior work on multi-party conversational AI, such as
the studies by Addlesee~\cite{addlesee2024ari} with the ARI robot and by
Axelsson and Skantze~\cite{axelsson2025furhat} with the Furhat platform,
invariably relies on cascaded pipelines to handle the additional complexity of
speaker identification, overlapping speech, and group-level turn management.
This gap represents both a limitation of the current state of the art and an
opportunity for future research.



% --- Turn-Taking in Spoken Dialogue  &  VAD / VAP ---
\input{chapters/section-turntaking-vadvap}

% --- Multi-Party CA, Research Gaps, Positioning ---
% ============================================================
%  \subsection: Multi-Party Conversational AI
%  \subsection: Research Gaps
%  \subsection: Positioning of AIutami
%  Parent: \section{State of the Art}
% ============================================================

\subsection{Multi-Party Conversational AI}
\label{subsec:multiparty-ca}

% ----------------------------------------------------------
\subsubsection{Why Multi-Party Is Different}
\label{sssec:multiparty-why}

Nearly all conversational agents in widespread use today---Siri, Alexa,
ChatGPT, GPT-4o voice mode---are engineered for \emph{dyadic} interaction: a
single human speaks to a single artificial interlocutor, and they take turns in
strict alternation.  Real human communication, however, is predominantly
group-based: meetings, classroom discussions, therapy sessions, and
brainstorming workshops all involve multiple participants whose contributions
interleave, overlap, and address different subsets of the group.

The transition from dyadic to multi-party conversation is not an incremental
extension; it constitutes a qualitative paradigm shift that introduces
fundamentally new computational and design
problems~\cite{gu2022mpc,zheng2022polyadic,traum2004issues}.  In a dyadic
exchange, the information flow is sequential: speaker~A produces an utterance,
speaker~B responds, and so forth in a predictable $A \!\to\! B \!\to\! A$
pattern.  In a multi-party setting, the information flow becomes graph-based:
any utterance can originate from any participant and can be directed at any
individual, any subset, or the entire group, giving rise to parallel
sub-conversations, shifting coalitions, and complex patterns of turn
allocation~\cite{gu2022mpc}.  These structural differences propagate into every
layer of a conversational AI system---from speaker identification and turn
management to response generation and addressee resolution---and demand
dedicated modelling strategies that have no counterpart in dyadic research.

The remainder of this subsection organises the multi-party literature along
three complementary axes: computational, design, and experiential.  It then
reviews the small number of systems that have been built and deployed, and
concludes with an analysis of the technical obstacles that explain why most
multi-party research has remained text-based.

% ----------------------------------------------------------
\subsubsection{Taxonomies of Multi-Party Conversational AI}
\label{sssec:multiparty-taxonomies}

Three taxonomies, originating from different research communities, provide
complementary lenses through which multi-party conversational AI can be
analysed.  Together they span the computational, design, and user-experience
dimensions of the problem.

\paragraph{WHO Says WHAT to WHOM---a computational taxonomy.}
Gu et al.~\cite{gu2022mpc} present a comprehensive IJCAI survey that
decomposes multi-party conversation (MPC) into three interrelated computational
problems.  The first, \emph{Speaker Modelling} (\textsc{who}), concerns the
identification and prediction of the current and next speaker: who is speaking,
and who will speak next?  Solutions range from Conditional Random Fields and
LSTMs to Transformer-based turn-taking predictors.  The second, \emph{Utterance
Modelling} (\textsc{what}), addresses the generation of the agent's
contribution: what should it say?  The literature offers both retrieval-based
methods (e.g., TopicBERT) and generation-based approaches (e.g., HeterMPC, and
more recently LLM prompting).  The third, \emph{Addressee Modelling}
(\textsc{whom}), tackles the problem of determining to whom an utterance is
directed: addressee recognition, dialogue disentanglement, and the
disambiguation of shared-channel messages.

Gu et al.\ observe that these three problems are typically studied in
isolation, yet they are deeply interdependent in practice: knowing \emph{who}
speaks conditions the interpretation of \emph{what} is said, which in turn
constrains the inference of \emph{whom} it addresses.  A critical limitation of
the survey, however, is that it covers exclusively \emph{text-based} MPC.  The
extension to speech introduces additional complexity layers---ASR errors,
overlapping speech, prosodic cues, and real-time latency
constraints---that the taxonomy does not address.

\paragraph{WHEN / WHAT / WHERE + SPECIFY / ACCESS / IMPLEMENT---a design taxonomy.}
Houde et al.~\cite{houde2025koala} propose an HCI-oriented framework,
developed at IUI~2025, for specifying and controlling the behaviour of an AI
agent participating in group conversations.  Their taxonomy is organised along
two orthogonal axes.

The first axis captures \emph{what to control}:
\begin{itemize}[nosep]
  \item \textsc{when}---the triggers, filters, and frequency that govern when
    the agent contributes (e.g., only when addressed, when a self-assessed
    relevance score exceeds a threshold, or at a fixed cadence);
  \item \textsc{what}---the content, style, and modality of the agent's
    contributions (e.g., conservative vs.\ creative ideas, formal vs.\ casual
    tone);
  \item \textsc{where}---the channel through which the agent communicates (a
    shared group channel, a threaded reply, or a private message).
\end{itemize}

The second axis captures \emph{how to control it}:
\begin{itemize}[nosep]
  \item \textsc{specify}---the mechanism for configuring the agent's behaviour
    (a UI panel, natural-language instructions, role-based presets);
  \item \textsc{access}---who is authorised to adjust the configuration (an
    administrator, any participant, or a democratic consensus);
  \item \textsc{implement}---how the configuration is enforced at the system
    level (system prompt, external logic module, or a hybrid of both).
\end{itemize}

This design taxonomy is complementary to Gu et al.'s computational
decomposition: where the latter describes the \emph{technical sub-problems} of
MPC, Houde et al.'s framework describes the \emph{design space} for governing
the agent's participation.

\paragraph{Four Challenges and Three Design Properties---a UX taxonomy.}
Zheng et al.~\cite{zheng2022polyadic} conduct a CHI literature review of 36
papers on polyadic (i.e., multi-party) human--AI interaction, distilling four
recurrent challenges:
\begin{enumerate}[nosep]
  \item \emph{Inefficient communication}---topic drift, lack of structure,
    difficulty reaching consensus;
  \item \emph{Lack of engagement}---unequal participation, passive members;
  \item \emph{Relationship-maintenance barriers}---insufficient emotional
    awareness, difficulty regulating group affect;
  \item \emph{Need for building connections}---ice-breaking, common-ground
    establishment, cross-cultural challenges.
\end{enumerate}

From these challenges they derive three design properties that a polyadic
conversational agent should exhibit.  The agent should be
\emph{Visible}---participants must be aware of its presence and role.  It
should be \emph{Ignorable}---its interventions must be non-intrusive; if they
are too frequent or too long, they become counterproductive.  And it should be
\emph{Accountable}---it must be clear to whom the agent is responding and who
controls its behaviour.

These three properties impose concrete constraints on system design.  In
text-based environments, visibility can be achieved through distinct usernames
or badges, ignorability through message threading, and accountability through
explicit @-mentions.  In speech-based environments, as we discuss in
Section~\ref{sssec:gaps}, all three properties take on fundamentally different
meanings: the agent's voice is inherently visible (it occupies a shared audio
channel), difficult to ignore (one cannot scroll past a spoken utterance), and
hard to make accountable without explicit turn-assignment mechanisms.

% ----------------------------------------------------------
\subsubsection{Existing Multi-Party Systems}
\label{sssec:multiparty-systems}

Despite the extensive theoretical groundwork laid by the taxonomies above, the
number of systems that actually implement a multi-party conversational agent
remains small.  This section reviews five such systems---spanning text, speech,
and multimodal modalities---and characterises each according to its interaction
modality, the role assumed by the AI, the turn-management mechanism, and the
degree of user control.

\paragraph{Koala (Houde et al., 2025).}
Koala~\cite{houde2025koala} is a Slack-based chatbot deployed in brainstorming
sessions at IBM, with three human participants and one AI agent.  Houde et al.\
evaluate two variants.  In the \emph{reactive} variant, Koala contributes only
when explicitly addressed (via an @-mention).  In the \emph{proactive} variant,
Koala employs a self-scoring LLM mechanism: after each human message, it
estimates the potential value of a contribution on a 0--100 scale and
autonomously posts if the score exceeds a configurable threshold.

The proactive variant generated 73\% of all ideas in the AI-augmented
conditions, yet only 33\% of the top-rated ideas were attributed to the agent.
More critically, 72\% of participants preferred the reactive variant, and the
proactive agent was described as a ``pedantic student who wouldn't create space
for others.''  A follow-up study, Koala~II, introduced a user-configurable
threshold: notably, no group reverted to the fully reactive mode, and the
perceived usefulness of the controls was rated 4.46 out of~5.  The central
finding is that agent proactivity should be neither binary nor static: users
want to adjust it dynamically during a session.

\paragraph{ARI Robot (Addlesee et al., 2024).}
Addlesee et al.~\cite{addlesee2024ari} deploy the humanoid robot ARI in a
hospital memory clinic, where it interacts with two human participants (a
patient and a companion) via speech.  The system follows a cascaded pipeline:
ASR transcription is followed by gaze-based addressee detection using an LLM
classifier, which achieves 85.4\% accuracy---compared with 53.4\% when using
text alone.  If the robot determines that it is not the intended addressee, it
takes no action.

Two additional contributions are noteworthy.  First, the system implements
\emph{incremental Clarification Requests} (iCR): when an utterance is detected
as incomplete---a frequent occurrence with patients experiencing cognitive
decline---the robot generates a natural follow-up question rather than
attempting to respond to a fragmentary input.  Second, a ``Jodie'' grounding
technique attributes information to a fictitious persona, forcing the LLM to
anchor its responses in the provided context and reducing hallucination by an
average of 10 percentage points across annotated responses.  A key limitation is
that the system has only been tested with two human participants and depends on
a gaze-tracking camera for addressee detection.

\paragraph{Furhat Robot (Axelsson and Skantze, 2025).}
Axelsson and Skantze~\cite{axelsson2025furhat} present a social robot built on
the Furhat platform that engages in open-ended, unscripted conversation with two
human participants.  The system's pipeline processes voice direction-of-arrival
signals, Azure STT transcription, speaker diarisation, and face recognition
before forwarding the result to GPT-3.5 for response generation; the output is
rendered via TTS and accompanied by gaze and head gestures.

In a parallel setting (two users interacting separately with the robot),
addressee accuracy reaches 92.6\%, but in a genuine group-conversation setting
it drops to 79.3\%.  The most striking result concerns \emph{audio-only}
addressee recognition for concurrent speakers: accuracy falls to just 26.8\%,
confirming that audio alone is insufficient for resolving addressee ambiguity in
multi-party speech.  Mean response latency is 1.35\,s, dominated by LLM
inference.  The system is limited to two human participants, operates in a
controlled laboratory environment, and assigns no moderating role to the robot.

\paragraph{Empathic Co-facilitator (Adikari et al., 2022).}
Adikari et al.~\cite{adikari2022empathic} propose a text-based AI
co-facilitator for online cancer-support groups (Cancer Chat Canada), analysing
approximately 120\,000 conversational turns.  The AI operates in the background
alongside a human therapist: it monitors group dynamics, detects emotions using
Plutchik's eight-emotion taxonomy via Word2Vec and GloVe embeddings, and models
emotion-state transitions through Markov chains.  When a negativity score
exceeds a threshold ($\lambda > 0.5$), the system either generates an empathic
response from pre-defined templates or alerts the human therapist.

The system demonstrates that an AI can play a valuable support role in
emotionally sensitive group settings.  However, it pre-dates the LLM era: its
natural-language capabilities are limited to template-based responses and
classical classifiers, and it operates exclusively in asynchronous text.

\paragraph{VIRTUAL HUMAN (Wahlster, 2006/2023).}
Wahlster~\cite{wahlster2023virtual} describes a historical precursor: a
multimodal sports-quiz system from 2006, featuring two human contestants and
three virtual agents (one moderator and two domain experts).  Dialogue
management relied on symbolic rules rather than statistical or neural models.
While the system demonstrates that the concept of an AI moderator in multi-party
conversation has a two-decade lineage, its pre-LLM architecture makes it
fundamentally different from contemporary approaches in both flexibility and
scalability.

\paragraph{Comparative synthesis.}
Table~\ref{tab:multiparty-comparison} provides a systematic comparison of all
five systems, together with AIutami, across ten dimensions that capture the
essential design choices identified by the taxonomies in
Section~\ref{sssec:multiparty-taxonomies}.

\begin{table}[ht]
  \centering
  \caption{Comparison of multi-party conversational AI systems.  AIutami is
    included for reference and discussed in detail in
    Section~\ref{subsec:positioning}.}
  \label{tab:multiparty-comparison}
  \footnotesize
  \setlength{\tabcolsep}{3.5pt}
  \begin{tabularx}{\textwidth}{@{} l X X X X X X @{}}
    \toprule
    \textbf{Dimension}
      & \textbf{Koala}
      & \textbf{ARI Robot}
      & \textbf{Furhat}
      & \textbf{Co-facilitator}
      & \textbf{Virtual Human}
      & \textbf{AIutami} \\
    \midrule
    Year
      & 2025 & 2024 & 2025 & 2022 & 2006 & 2024 \\
    Modality
      & Text (Slack)
      & Speech + gaze
      & Speech + gaze
      & Text (chat)
      & Multimodal
      & Speech (WebRTC) \\
    AI role
      & Participant
      & Receptionist
      & Conv.\ partner
      & Co-facilitator
      & Moderator
      & Moderator \\
    Participants
      & 3\,H + 1\,AI
      & 2\,H + 1\,robot
      & 2\,H + 1\,robot
      & $N$\,H + AI + therapist
      & 2\,H + 3\,agents
      & $N$\,H + 1\,AI \\
    \textsc{when}
      & Self-scoring (0--100)
      & Addressee (gaze)
      & Gaze + silence
      & Negativity $\lambda$
      & Symbolic rules
      & Trigger-based \\
    User control
      & In-session panel
      & None
      & None
      & None (therapist)
      & None
      & Session config \\
    LLM
      & Llama\,2/3
      & Vicuna-13b
      & GPT-3.5
      & Pre-LLM
      & Pre-LLM
      & Azure OpenAI \\
    Turn-taking
      & Implicit (text)
      & Addressee det.
      & Gaze-based
      & N/A (async)
      & Multimodal rules
      & Explicit + 8\,s resv. \\
    Context
      & Brainstorming
      & Memory clinic
      & Open-ended
      & Oncology support
      & Sports quiz
      & Multi-context \\
    Deployed
      & Prototype
      & Hospital
      & Lab prototype
      & Prototype
      & Historical
      & Web (production) \\
    \bottomrule
  \end{tabularx}
\end{table}

% ----------------------------------------------------------
\subsubsection{Technical Challenges of Multi-Party Speech}
\label{sssec:multiparty-technical}

The paucity of speech-based multi-party systems is not accidental: it reflects
the severity of the technical challenges that speech introduces relative to
text.  Addlesee et al.~\cite{addlesee2020asr} provide a systematic evaluation
of commercial ASR and speaker-diarisation services in both dyadic and
multi-party settings, revealing performance degradation that is often
prohibitive.  Table~\ref{tab:asr-challenges} summarises their findings.

\begin{table}[ht]
  \centering
  \caption{Technical challenges of ASR in multi-party spoken interaction,
    adapted from Addlesee et al.~\cite{addlesee2020asr}.}
  \label{tab:asr-challenges}
  \small
  \begin{tabularx}{\textwidth}{@{} l X l @{}}
    \toprule
    \textbf{Challenge} & \textbf{Description} & \textbf{Measured Impact} \\
    \midrule
    Speaker diarisation
      & Determining who is speaking at each moment
      & DER 49--67\% (multi-party) \\
    Overlapping speech
      & Concurrent speakers degrade recognition
      & ${\sim}$20\% WER degradation \\
    Incremental ASR
      & Real-time partial results vs.\ batch
      & 33\% WER (incr.)\ vs.\ 5\% (batch) \\
    Disfluencies
      & Filled pauses, self-corrections, edit terms
      & Variable across services \\
    \bottomrule
  \end{tabularx}
\end{table}

Speaker diarisation---the task of attributing each audio segment to the correct
speaker---exhibits Diarisation Error Rates between 49\% and 67\% in multi-party
settings on commercial platforms, with Google's service failing entirely when
the number of speakers exceeds three.  Overlapping speech, which occurs
naturally whenever speakers compete for the floor or produce backchannels,
degrades Word Error Rates by approximately 20\%.  Incremental (streaming) ASR,
which is necessary for responsive turn-taking, yields a sixfold increase in WER
compared with batch transcription (33\% vs.\ 5\%).  These compound difficulties
explain why the vast majority of multi-party conversational AI research---as
surveyed by Gu et al.~\cite{gu2022mpc}, Zheng et
al.~\cite{zheng2022polyadic}, and Houde et al.~\cite{houde2025koala}---has
operated exclusively in the text modality.


% ============================================================
\subsection{Research Gaps}
\label{subsec:gaps}

The preceding sections have traced three research threads---speech-based
architectures, turn-taking models, and multi-party conversational
AI---and have shown how each has matured largely in isolation from the others.
This subsection consolidates the gaps that emerge at their intersection,
identifying eight open problems that delineate opportunities for future work.

% ----------------------------------------------------------
\subsubsection{Gap~1: Multi-Party Speech-Based AI---the Nearly Empty Quadrant}
\label{sssec:gap1}

The most prominent gap in the literature is the near-absence of systems that
combine speech-based interaction with multi-party group dynamics and an active
AI participant.  As shown in Table~\ref{tab:multiparty-comparison}, only the
ARI robot~\cite{addlesee2024ari} and the Furhat
robot~\cite{axelsson2025furhat} operate in speech, and both are restricted to
exactly two human participants plus one robot, with a hard dependency on gaze
tracking for addressee resolution.  No existing system simultaneously satisfies
all of the following criteria: speech-based input and output, support for $N$
human participants, an active AI moderator role, LLM-based reasoning, and
audio-only operation without gaze or video.  The MPCA
Survey~\cite{mpcasurvey2025} corroborates this finding indirectly: it does not
systematically distinguish between text-based and speech-based multi-party
agents, signalling that the speech dimension is underrepresented in the field.

% ----------------------------------------------------------
\subsubsection{Gap~2: End-to-End and Half-Cascade Architectures for Multi-Party}
\label{sssec:gap2}

All recent speech-to-speech models---Moshi~\cite{defossez2024moshi},
GPT-4o, LLaMA-Omni~\cite{fang2024llamaomni},
Freeze-Omni~\cite{wang2024freezeomni}, Qwen3-Omni---are
designed exclusively for dyadic interaction.  No published work investigates
how these architectures could be adapted to the multi-party setting, where the
system must handle $N$ concurrent audio streams, perform real-time speaker
identification, manage $N$-way turn allocation, and decide \emph{when} and
\emph{to whom} to respond.  The only architecture with demonstrated multi-party
support remains the cascaded STT--LLM--TTS pipeline, as employed by both
Addlesee et al.~\cite{addlesee2024ari} and AIutami.

% ----------------------------------------------------------
\subsubsection{Gap~3: Design Taxonomy for Multi-Party Speech}
\label{sssec:gap3}

Houde et al.'s design taxonomy~\cite{houde2025koala} was developed and
validated exclusively on Slack, a text-based platform.  Transposing it to
speech fundamentally alters the design space in at least four ways.  First,
\emph{latency matters acutely}: in text, a response delivered after five
seconds is unremarkable; in speech, a five-second silence is socially
disruptive, and the agent must decide whether to intervene in under one second.
Second, \emph{voice is sequential and non-skippable}: a long text message can
be scanned or scrolled past, but a long spoken utterance physically occupies
the shared audio channel and cannot be fast-forwarded.  Third, \emph{there are
no threads in speech}: the \textsc{where} dimension collapses to a single
shared channel (or, at best, a parallel private channel), eliminating the
threaded-reply affordance that text platforms provide.  Fourth, \emph{control
interfaces are incompatible with ongoing speech}: a participant cannot click a
configuration panel while speaking.  DialogLab~\cite{hu2025dialoglab},
a recent authoring tool from Google presented at UIST~2025, enables the design
and simulation of group human--AI conversations but does not propose a design
taxonomy specific to speech, nor does it address the latency and sequentiality
constraints of the vocal modality.

% ----------------------------------------------------------
\subsubsection{Gap~4: Audio-Only Addressee Detection in Groups with $N > 2$}
\label{sssec:gap4}

Addressee detection---determining to whom an utterance is directed---is a
prerequisite for appropriate agent behaviour in multi-party settings.  Current
results are discouraging for audio-only scenarios.  Addlesee et
al.~\cite{addlesee2024ari} achieve 85.4\% accuracy using a combination of gaze
and text features, but only with two human participants; text alone yields
53.4\%.  Axelsson and Skantze~\cite{axelsson2025furhat} report 92.6\% accuracy
with gaze features, but when restricted to audio-only input with concurrent
speakers, accuracy collapses to 26.8\%.  Without gaze or video, and with more
than two speakers, the problem remains essentially open.  Lee and
Deng~\cite{lee2024endofturn}, Best Paper Runner-up at ICMI~2024, address
end-of-turn prediction for three-party conversations, but their approach
requires an eight-camera motion-capture setup---far removed from the
constraints of a typical web-based deployment.

% ----------------------------------------------------------
\subsubsection{Gap~5: Dynamic Agent Control in Speech Contexts}
\label{sssec:gap5}

Houde et al.~\cite{houde2025koala} demonstrate that users strongly prefer the
ability to adjust agent behaviour during a session rather than only before it
(rated 4.46/5 for perceived usefulness of in-session controls).  In text-based
environments, this is straightforward: a side panel, a slider, or an in-chat
command.  In speech, however, the control and conversational channels are
conflated.  A voice command such as ``moderator, speak less'' is ambiguous:
it could be interpreted as an instruction to the agent or as a conversational
utterance directed at another participant.  Alternative modalities---gestural
input (if video is available), a parallel web interface, or keyword-triggered
control commands---have not been explored in the literature.

% ----------------------------------------------------------
\subsubsection{Gap~6: Incomplete Utterance Handling in Multi-Party Speech}
\label{sssec:gap6}

Addlesee et al.~\cite{addlesee2024ari} introduce incremental Clarification
Requests (iCR) to handle incomplete utterances from patients with dementia.  In
generic multi-party speech, the problem is amplified.  ASR endpointing
algorithms may segment an utterance mid-sentence, particularly when the
speaker pauses to think.  With $N$ speakers, interruptions and false endpoints
occur more frequently, and the downstream LLM receives truncated input that
can trigger incoherent or irrelevant responses.  No multi-party speech system
currently implements recovery strategies---such as iCR, utterance-continuation
detection, or endpointing rollback---for the general case.

% ----------------------------------------------------------
\subsubsection{Gap~7: Benchmarks for Multi-Party Speech-Based CA}
\label{sssec:gap7}

The field lacks a shared evaluation infrastructure for multi-party speech-based
conversational agents.  There is no standard dataset, no common benchmark, and
no agreed-upon evaluation protocol that spans both technical metrics (latency,
WER, diarisation error rate) and user-experience measures (discussion quality,
participation balance, perceived group climate).  Zheng et
al.~\cite{zheng2022polyadic} propose UX evaluation dimensions for polyadic
agents, but only for text.  Addlesee et al.~\cite{addlesee2020asr} provide ASR
metrics for multi-party speech, but without downstream LLM evaluation.  Recent
spoken-dialogue benchmarks---SD-Eval and Full-Duplex-Bench---are exclusively
dyadic.  The MPCA Survey~\cite{mpcasurvey2025} explicitly confirms this gap,
noting that ``current MPCA evaluation benchmarks have three limitations:
(i)~each dataset focuses on a single specific skill of MPCAs; (ii)~lack of
benchmarks for multi-modal settings; (iii)~lack of realistic
metrics/simulations for evaluation.''

% ----------------------------------------------------------
\subsubsection{Gap~8: Amplified Proactivity Risk in Speech}
\label{sssec:gap8}

Houde et al.~\cite{houde2025koala} show that even in text, a proactive agent
was perceived as a ``pedantic student who wouldn't create space for others.''
In speech, this risk is amplified by three factors.  First, the AI's voice
occupies a shared audio channel---unlike a text message, it cannot be scrolled
past or collapsed.  Second, a long intervention physically blocks the
conversation: all human participants must wait until the agent finishes
speaking, an effect known as \emph{production blocking} in the group-creativity
literature.  Third, the sequential nature of speech means that even a modestly
verbose intervention consumes disproportionate floor time.  The balance between
active moderation and conversational non-intrusiveness has not been empirically
studied in speech-based multi-party settings.


% ============================================================
\subsection{Positioning of AIutami}
\label{subsec:positioning}

% ----------------------------------------------------------
\subsubsection{The Research Landscape as a 2$\times$2 Matrix}
\label{sssec:matrix}

The preceding review reveals a pronounced asymmetry in the literature.
Table~\ref{tab:positioning-matrix} organises the surveyed systems along two
dimensions---interaction modality (text vs.\ speech) and conversation topology
(dyadic vs.\ multi-party)---and exposes the quadrant in which AIutami is
situated.

\begin{table}[ht]
  \centering
  \caption{Positioning of conversational AI systems along two dimensions.  The
    speech--multi-party quadrant is extremely sparse.}
  \label{tab:positioning-matrix}
  \small
  \begin{tabularx}{\textwidth}{@{} l >{\raggedright\arraybackslash}X >{\raggedright\arraybackslash}X @{}}
    \toprule
    & \textbf{Text-Based} & \textbf{Speech-Based} \\
    \midrule
    \textbf{Dyadic} \newline (1 human + 1 AI)
      & ChatGPT, Siri, Alexa \newline (extensive literature)
      & GPT-4o, Moshi, LLaMA-Omni, Freeze-Omni \newline (growing literature) \\
    \addlinespace
    \textbf{Multi-party} \newline ($N$ humans + AI)
      & Koala~\cite{houde2025koala}, Adikari~\cite{adikari2022empathic},
        DialogLab~\cite{hu2025dialoglab} \newline (moderate literature)
      & \textbf{AIutami}, Addlesee~\cite{addlesee2024ari}\,(2+1, gaze),
        Furhat~\cite{axelsson2025furhat}\,(2+1, gaze) \newline
        \emph{(extremely sparse)} \\
    \bottomrule
  \end{tabularx}
\end{table}

The text--dyadic and speech--dyadic quadrants enjoy substantial coverage, and
the text--multi-party quadrant has seen growing activity in recent years.  The
speech--multi-party quadrant, by contrast, contains only two prior systems, both
of which are limited to two human participants, depend on gaze tracking, and do
not assign a moderating role to the AI.  To the best of our knowledge, few existing systems combine all of the
following properties: (i)~speech-based input and output, (ii)~support for
$N$ concurrent human participants, (iii)~an active AI moderator role,
(iv)~audio-only operation with no dependency on gaze or video, and
(v)~LLM-based reasoning.  AIutami is an attempt to address this
combination within a single platform.

% ----------------------------------------------------------
\subsubsection{How AIutami Addresses Known Challenges}
\label{sssec:aiutami-solutions}

AIutami's architecture incorporates design decisions that directly address
several of the technical and design challenges identified in the literature.
Table~\ref{tab:aiutami-solutions} maps each challenge to the corresponding
architectural mechanism.

\begin{table}[ht]
  \centering
  \caption{How AIutami addresses challenges identified in the literature.}
  \label{tab:aiutami-solutions}
  \small
  \begin{tabularx}{\textwidth}{@{} l X l @{}}
    \toprule
    \textbf{Challenge (source)} & \textbf{AIutami's approach} & \textbf{Module} \\
    \midrule
    Unreliable speaker diarisation \newline\cite{addlesee2020asr}
      & Separate WebRTC audio streams per participant eliminate the need for diarisation entirely
      & WebRTC \\
    \addlinespace
    Overlapping speech degrades ASR \newline\cite{addlesee2020asr}
      & ASR gating: transcription is activated only for the current speaker, as determined by the turn-management subsystem, preventing echo and cross-talk from contaminating the input
      & ASR \\
    \addlinespace
    Addressee detection without gaze \newline\cite{addlesee2024ari}
      & Structural bypass via explicit turn-taking: the moderator addresses the group, and participants request the floor through a dedicated control, removing addressee ambiguity
      & Turns \\
    \addlinespace
    Disruptive proactivity \newline\cite{houde2025koala}
      & 8-second reservation window for the next speaker after the current turn ends, combined with conditional trigger logic that limits unprompted AI interventions
      & Turns + Moderation \\
    \addlinespace
    Unequal engagement \newline\cite{zheng2022polyadic}
      & Balanced turn-taking mechanism: any participant can request a turn at any time, and the moderator can explicitly invite quieter members
      & Turns \\
    \addlinespace
    Inefficient communication \newline\cite{zheng2022polyadic}
      & Evolving summary maintained by the AI; the moderator manages discussion flow, redirects off-topic contributions, and synthesises emerging consensus
      & Moderation \\
    \bottomrule
  \end{tabularx}
\end{table}

Three architectural choices merit additional discussion.  First, the use of
\emph{per-participant WebRTC audio streams} is a deliberate departure from the
single-channel audio that characterises both ARI and Furhat.  By receiving a
dedicated stream from each participant, AIutami sidesteps the speaker
diarisation problem entirely---a problem that, as Addlesee et
al.~\cite{addlesee2020asr} demonstrate, remains largely unsolved in real-time
multi-party settings (DER 49--67\%).

Second, the \emph{ASR gating} mechanism---whereby transcription is activated
only for the participant who currently holds the floor---mitigates two problems
at once.  It prevents overlapping speech from corrupting the transcript (thereby
avoiding the ${\sim}$20\% WER degradation reported by Addlesee et
al.~\cite{addlesee2020asr}), and it eliminates echo loops in which the agent's
own TTS output is inadvertently re-transcribed.

Third, the \emph{explicit turn-taking} model with an 8-second reservation
window provides a structural solution to both the addressee-detection and the
proactivity challenges.  Rather than attempting to infer the addressee from
prosodic or textual cues---a task that achieves at most 26.8\% accuracy in
audio-only multi-party settings~\cite{axelsson2025furhat}---AIutami eliminates
the ambiguity by design: the floor is assigned to one speaker at a time, and
the moderator's interventions are directed at the group as a whole.  The
reservation window ensures that the AI does not immediately seize the floor
after a human finishes speaking, providing a temporal buffer that reduces the
``production blocking'' effect observed by Houde et
al.~\cite{houde2025koala}.

% ----------------------------------------------------------
\subsubsection{Limitations and Open Directions}
\label{sssec:aiutami-limitations}

The current design of AIutami leaves several gaps unaddressed, each of which
represents a potential direction for future work.
Table~\ref{tab:aiutami-limitations} summarises these limitations alongside the
literature that motivates them.

\begin{table}[ht]
  \centering
  \caption{Challenges not addressed by AIutami and their corresponding
    literature references.}
  \label{tab:aiutami-limitations}
  \small
  \begin{tabularx}{\textwidth}{@{} l X l @{}}
    \toprule
    \textbf{Open challenge} & \textbf{Potential extension} & \textbf{Reference} \\
    \midrule
    Clarification requests for incomplete utterances
      & Detect truncated utterances at the ASR endpointing stage and generate
        natural follow-up questions before forwarding to the LLM
      & \cite{addlesee2024ari} \\
    \addlinespace
    Dynamic agent control during session
      & Voice commands with keyword detection, a parallel web interface, or
        gestural input to adjust moderator behaviour in real time
      & \cite{houde2025koala} \\
    \addlinespace
    Explicit anti-hallucination grounding
      & Adopt context-anchoring techniques (e.g., the ``Jodie'' persona) to
        force the LLM to ground responses in session-specific information
      & \cite{addlesee2024ari} \\
    \addlinespace
    Emotion-aware moderation
      & Integrate sentiment analysis from transcribed text or prosodic features
        into the moderator's decision logic, enabling affect-sensitive
        interventions
      & \cite{adikari2022empathic} \\
    \addlinespace
    VAP-informed turn management
      & Incorporate Voice Activity Projection~\cite{ekstedt2022vap} to
        anticipate turn transitions, reducing perceived latency and enabling
        smoother floor handovers
      & \cite{ekstedt2022vap} \\
    \bottomrule
  \end{tabularx}
\end{table}

These limitations are not incidental: they correspond directly to
Gaps~4--8 identified in Section~\ref{subsec:gaps} and represent the most
promising avenues for extending the current system.  The explicit turn-taking
model, while effective, sacrifices conversational naturalness; future work might
relax it by integrating predictive turn-taking signals.  The absence of
emotion-aware moderation means that the system currently optimises for
\emph{procedural} fairness (equal floor time) without attending to
\emph{affective} dynamics (group mood, individual distress).  And the lack of
in-session control mechanisms means that, unlike Koala~II's configurable
threshold, AIutami's moderator behaviour is fixed for the duration of a
session.

Together, the solutions and limitations outlined in this section position
AIutami as an early attempt to populate the largely empty
speech--multi-party quadrant---one that addresses several critical technical
barriers (diarisation, overlapping speech, addressee ambiguity, proactivity
risk) while leaving a well-defined set of open challenges for future
investigation.

